\documentclass[14pt]{extarticle}
\usepackage{fontspec}
\usepackage{polyglossia}
\usepackage{amsmath}
\usepackage{amssymb}
\usepackage{xcolor}
\usepackage{fancyhdr}
\usepackage{graphicx}
\usepackage{listings}
\usepackage[most]{tcolorbox}
% \usepackage{setspace}
\usepackage{pifont}
\usepackage{enumitem}
\usepackage{titlesec}
\usepackage[bottom]{footmisc}
\usepackage{titling}
\usepackage{minted}
\usepackage{etoolbox}
\usepackage{array}
\usepackage{longtable} % page-breakable tables (must load before polyglossia/bidi RTL setup)
\usepackage{extsizes}
\usepackage{float}
\usepackage{booktabs}
\usepackage{hyperref}
\usepackage[onehalfspacing]{setspace}
\newif\ifwhiteblack
\ifdefined\setwhiteblack
  \setwhiteblack
\fi

\ifwhiteblack
  \usepackage[a4paper, landscape, top=1.25cm, bottom=1.25cm, left=1.25cm, right=1.25cm, headheight=1.5cm, headsep=0.5cm]{geometry}
\else
  \usepackage[a4paper, top=2.7cm, bottom=2.7cm, left=2.5cm, right=2.5cm, headheight=1.5cm, headsep=0.5cm]{geometry}
\fi
\usepackage{tikz}
\usepackage{datetime2}

\usetikzlibrary{automata, positioning, arrows.meta, arrows}
\providecommand{\LTcaptype}{table}

\newfontfamily\emoji{Segoe UI Emoji}

\pagestyle{fancy}

\setmainlanguage[numerals=western]{arabic}
\setotherlanguage{english}
\newfontfamily\arabicfont[Script=Arabic]{Calibri}
\newfontfamily\arabicfonttt[Script=Arabic]{Courier New}



\ifwhiteblack
\lstset{
  language=[Sharp]C,
  numbers=left,
  stepnumber=1,
  numbersep=8pt,
  frame=single,
  basicstyle=\ttfamily\small,
  keywordstyle=\bfseries,
  stringstyle=\ttfamily,
  commentstyle=\itshape
}
\else
\lstset{
  language=[Sharp]C,
  numbers=left,
  stepnumber=1,
  numbersep=8pt,
  frame=single,
  basicstyle=\ttfamily\small,
  keywordstyle=\color{blue},
  stringstyle=\color{red},
  commentstyle=\color{green!50!black}
}
\fi

\newif\ifdetailed
\ifdefined\setdetailed
  \setdetailed
\fi

\newif\ifwithsols
\ifdefined\setwithsols
  \setwithsols
\fi

\ifwhiteblack
  \AtBeginDocument{
    \LARGE
    \usemintedstyle{bw}
  }
\fi

% Default minted look: compact code that wraps long lines.
% Not applied in black-and-white landscape mode, which already uses \LARGE.
\ifwhiteblack\else
  \setminted{fontsize=\small, breaklines=true}
\fi

% unified tcolorboxes for articles
\ifwhiteblack
\tcbset{colback=white, colframe=black, fonttitle=\bfseries, boxrule=1pt}
\newtcolorbox{boxDef}[1][]{title={{\emoji📘} تعريف\ifx\\#1\\\else ~#1\fi :}}
\newtcolorbox{boxExercise}[1][]{title={{\emoji🧩} تمرين\ifx\\#1\\\else ~#1\fi :}}
\newtcolorbox{boxExample}[1][]{title={{\emoji📝} مثال\ifx\\#1\\\else ~#1\fi :}}
\newtcolorbox{boxNote}[1][]{title={{\emoji✨} ملاحظة\ifx\\#1\\\else ~#1\fi :}}
\newtcolorbox{boxAttention}[1][]{title={{\emoji🔔} تنبيه\ifx\\#1\\\else ~#1\fi :}}
\newtcolorbox{boxWarning}[1][]{title={{\emoji⚡} ملاحظة هامة\ifx\\#1\\\else ~#1\fi :}}
\newtcolorbox{boxSolution}[1][]{title={{\emoji✅} حل\ifx\\#1\\\else ~#1\fi :}}
\newtcolorbox{boxSymbol}[1][]{title={{\emoji🔣} رمز\ifx\\#1\\\else ~#1\fi :}}
\newtcolorbox{boxHint}[1][]{title={{\emoji💡} تلميح\ifx\\#1\\\else ~#1\fi :}}
\else
\tcbset{colback=white, colframe=black, fonttitle=\bfseries, boxrule=0.8pt}
\newtcolorbox{boxDef}[1][]{colback=blue!5!white,colframe=blue!75!black,
  title={{\emoji📘} تعريف\ifx\\#1\\\else ~#1\fi :}}
\newtcolorbox{boxExercise}[1][]{colback=cyan!5!white,colframe=cyan!70!black,
  title={{\emoji🧩} تمرين\ifx\\#1\\\else ~#1\fi :}}
\newtcolorbox{boxExample}[1][]{colback=yellow!5!white,colframe=orange!90!black,
  title={{\emoji📝} مثال\ifx\\#1\\\else ~#1\fi :}}
\newtcolorbox{boxNote}[1][]{colback=gray!10!white,colframe=black,
  title={{\emoji✨} ملاحظة\ifx\\#1\\\else ~#1\fi :}}
\newtcolorbox{boxAttention}[1][]{colback=magenta!10!white,colframe=magenta!80!black,
  title={{\emoji🔔} تنبيه\ifx\\#1\\\else ~#1\fi :}}
\newtcolorbox{boxWarning}[1][]{colback=red!5!white,colframe=red!75!black,
  title={{\emoji⚡} ملاحظة هامة\ifx\\#1\\\else ~#1\fi :}}
\newtcolorbox{boxSolution}[1][]{colback=green!5!white,colframe=green!60!black,
  title={{\emoji✅} حل\ifx\\#1\\\else ~#1\fi :}}
\newtcolorbox{boxSymbol}[1][]{colback=purple!5!white,colframe=purple!70!black,
  title={{\emoji🔣} رمز\ifx\\#1\\\else ~#1\fi :}}
\newtcolorbox{boxHint}[1][]{colback=teal!5!white,colframe=teal!60!black,
  title={{\emoji💡} تلميح\ifx\\#1\\\else ~#1\fi :}}
\fi


\ifwhiteblack
\tcbset{simplecode/.style={ colback=white, colframe=black, boxrule=1pt, arc=2pt, left=4pt,right=4pt,top=4pt,bottom=4pt}}
\else
\tcbset{simplecode/.style={ colback=gray!5, colframe=black!50, boxrule=0.4pt, arc=2pt, left=4pt,right=4pt,top=4pt,bottom=4pt}}
\fi
% Compact code boxes: tighter line spacing (the preamble uses onehalfspacing) and small vertical gaps.
% Must come AFTER the simplecode definition above.
\tcbset{simplecode/.append style={fontupper=\setstretch{0.95}, top=4pt, bottom=4pt, before skip=6pt, after skip=6pt}}
\newenvironment{boxCode}{\begin{tcolorbox}[simplecode]}{\end{tcolorbox}}

% Code cell used inside tables (keeps code left-to-right inside RTL tables).
% Use inside a fixed-width p{...} column: \begin{pycell}\begin{minted}...\end{minted}\end{pycell}
\newtcolorbox{pycell}{frame empty,colback=white,boxsep=0pt,left=2pt,right=2pt,top=2pt,bottom=2pt,width=\linewidth}

\newcolumntype{C}[1]{>{\centering\arraybackslash}p{#1}}

\makeatletter
\DTMnewdatestyle{dotted}{%
  \renewcommand{\DTMdisplaydate}[4]{\two@digits{##3}.\two@digits{##2}.##1}%
  \renewcommand{\DTMDisplaydate}{\DTMdisplaydate}%
}
\makeatother
\DTMsetdatestyle{dotted}

% redefine spaces after titles
\makeatletter
\renewcommand{\@maketitle}{%
  \begin{center}
    {\huge \bfseries \@title \par}%
    \vskip 0.2em % space between title and author
    {\large \@author \par}%
    % \vskip 0.2em % space between author and date
    % {\normalsize \@date \par}%
    {\small \hfill \begin{english}\DTMtoday\end{english} \par}%
  \end{center}
}
\makeatother

\fancyhf{} % clear default
\fancypagestyle{plain}{
  \fancyhf{}
  \fancyhead[L]{مدرسة التسامح الشاملة}
  % \fancyhead[L]{\includegraphics[height=1cm]{../../../images/logoTasamoh.png}}
  \fancyhead[R]{الأستاذ محمود اغبارية}
  \fancyfoot[C]{\thepage}
}

\fancyhead[L]{مدرسة التسامح الشاملة}
\fancyhead[R]{الأستاذ محمود اغبارية}
\fancyfoot[C]{\thepage}

\setcounter{tocdepth}{3} % only section subsection and subsubsection in TOC

\newcommand{\importpdfpage}[2]{%
    \thispagestyle{empty}
    \newgeometry{left=0cm,right=0cm,top=0cm,bottom=0cm}
    \includegraphics[page=#2,width=0.95\paperwidth,height=0.95\paperheight,keepaspectratio]{#1}
    \restoregeometry
}

\newcommand{\insertFullImg}[1]{%
    \noindent
    \makebox[\textwidth][c]{
        \includegraphics[width=0.9\paperwidth,keepaspectratio]{#1}%
    }%
}%


\ifwithsols
\title{حل امتحان مفمار في علم الحاسوب - بايثون (أ) $2025$}
\else
\title{امتحان مفمار في علم الحاسوب - بايثون (أ) $2025$}
\fi

\begin{document}

\maketitle
\thispagestyle{fancy}

\ifwithsols
\else
\begin{boxCode}
\begin{itemize}[nosep]
    \item \textbf{الفصل الأول (إجباري - $60$ علامة):} يجب الإجابة عن \textbf{جميع} الأسئلة $1$–$5$.
    \item \textbf{الفصل الثاني (اختياري - $40$ علامة):} يجب الإجابة عن \textbf{سؤالين} من بين الأسئلة $6$–$9$.
    \item \textbf{مدة الامتحان:} $120$ دقيقة.
    \item \textbf{مواد مساعدة مسموحة:} كل مادة مساعدة شخصية مكتوبة أو مطبوعة، باستثناء حاسبة أو حاسوب قابل للبرمجة.
    \item اقرؤوا أسئلة الامتحان بتمعّن وأجيبوا عنها بانتباه في المكان المخصّص لذلك.
\end{itemize}
\end{boxCode}
\fi

\section*{الفصل الأول — إجباري ($60$ علامة)}
في هذا الفصل خمسة أسئلة. يجب الإجابة عن \textbf{جميع} الأسئلة $1$–$5$.

\begin{enumerate}[label=\textbf{السؤال \arabic*:}, wide=0pt, itemsep=2.2em]

% ---------------------------------------------------------------
\item ($12$ علامة، $3$ علامات لكل بند) أمامكم جدول فيه أربعة مقاطع برنامج بلغة بايثون. لكل مقطع برنامج في الجدول، اكتبوا ما هو المُخرج الناتج في العمود الفارغ.

{\renewcommand{\arraystretch}{1.4}
\begin{longtable}{|C{1.5cm}|p{9cm}|C{4cm}|}
\hline
\textbf{البند} & \centering\textbf{كود بلغة بايثون} & \textbf{المُخرج الناتج} \\
\hline
$1$ &
\begin{pycell}
\begin{english}
\begin{minted}{python}
a = 7
b = a + 3
a = b - 5
print("a = ", a, " b = ", b)
\end{minted}
\end{english}
\end{pycell}
& \ifwithsols \begin{english}\small\texttt{a = 5 b = 10}\end{english} \fi \\
\hline
$2$ &
\begin{pycell}
\begin{english}
\begin{minted}{python}
value = 15
result = value // 2
value = result + 3
print("value = ", value, " result = ", result)
\end{minted}
\end{english}
\end{pycell}
& \ifwithsols \begin{english}\small\texttt{value = 10 result = 7}\end{english} \fi \\
\hline
$3$ &
\begin{pycell}
\begin{english}
\begin{minted}{python}
num = 7890
num1 = (num // 100) % 10
num2 = num % 100
print("num1 = ", num1, " num2 = ", num2)
\end{minted}
\end{english}
\end{pycell}
& \ifwithsols \begin{english}\small\texttt{num1 = 8 num2 = 90}\end{english} \fi \\
\hline
$4$ &
\begin{pycell}
\begin{english}
\begin{minted}{python}
num = 47
x = num // 10 + 10 * (num % 10)
print("num = ", num, " x = ", x)
\end{minted}
\end{english}
\end{pycell}
& \ifwithsols \begin{english}\small\texttt{num = 47 x = 74}\end{english} \fi \\
\hline
\end{longtable}}
\filbreak

% ---------------------------------------------------------------
\item ($12$ علامة، $3$ علامات لكل بند) في الجدول الذي أمامكم أربعة مقاطع برنامج جزئية والمُخرج الذي سينتج من تشغيل مقطع البرنامج الكامل. كل مقطع برنامج يحتوي على حلقة.
أكملوا التعليمات الناقصة في كل مقطع برنامج للحصول على المُخرج المعطى. استعينوا بالمثال في أعلى الجدول.

{\renewcommand{\arraystretch}{1.4}
\begin{longtable}{|C{1.5cm}|p{9cm}|C{4cm}|}
\hline
\textbf{البند} & \centering\textbf{كود بلغة بايثون} & \textbf{المُخرج الناتج} \\
\hline
مثال &
\begin{pycell}
\begin{english}
\begin{minted}{python}
for item in range(2, 7, 2):
    print(item, end=" ")
\end{minted}
\end{english}
\end{pycell}
& \begin{english}\texttt{2 4 6}\end{english} \\
\hline
$1$ &
\begin{pycell}
\begin{english}
\begin{minted}{python}
for item in range(2, _____):
    print(item, end=" ")
\end{minted}
\end{english}
\end{pycell}
& \begin{english}\texttt{2 3}\end{english} \\
\hline
$2$ &
\begin{pycell}
\begin{english}
\begin{minted}{python}
for item in range(10):
    if item % 3 == _________:
        print(item, end=" ")
\end{minted}
\end{english}
\end{pycell}
& \begin{english}\texttt{1 4 7}\end{english} \\
\hline
$3$ &
\begin{pycell}
\begin{english}
\begin{minted}{python}
item = __________
while item < 10:
    print(item, end=" ")
    item = ___________
\end{minted}
\end{english}
\end{pycell}
& \begin{english}\texttt{2 5 8}\end{english} \\
\hline
$4$ &
\begin{pycell}
\begin{english}
\begin{minted}{python}
item = 3
while ______ > ______:
    print(item * 5)
    item = item - 1
\end{minted}
\end{english}
\end{pycell}
& \begin{english}\texttt{15}\newline\texttt{10}\newline\texttt{5}\end{english} \\
\hline
\end{longtable}}

\ifwithsols
\begin{boxSolution}
\begin{enumerate}[label=\arabic*.]
\item \textenglish{\texttt{for item in range(2, 4):}}
\item \textenglish{\texttt{if item \% 3 == 1:}}
\item التهيئة \textenglish{\texttt{item = 2}}، والتحديث داخل الحلقة \textenglish{\texttt{item = item + 3}}
\item الشرط \textenglish{\texttt{while item > 0:}}
\end{enumerate}
\end{boxSolution}
\fi
\filbreak

% ---------------------------------------------------------------
\item ($12$ علامة) أمامكم عمودان من مقاطع البرنامج التي تُنفَّذ فيها حلقة. يجب الملاءمة بين أزواج الحلقات المتكافئة (بحيث ينتج في مقطعي الحلقة نفس المُخرج تمامًا) وكتابتها في الجدول أدناه.

{\renewcommand{\arraystretch}{1.3}
\begin{longtable}{|C{0.9cm}|p{5.8cm}|C{0.9cm}|p{5.8cm}|}
\hline
\multicolumn{2}{|c|}{\textbf{العمود أ}} & \multicolumn{2}{c|}{\textbf{العمود ب}} \\
\hline
$1$ &
\begin{pycell}
\begin{english}
\begin{minted}{python}
i = 5
while i <= 30:
    print(i)
    i = i + 5
\end{minted}
\end{english}
\end{pycell}
& $A$ &
\begin{pycell}
\begin{english}
\begin{minted}{python}
for i in range(16, 2, -3):
    print(i)
\end{minted}
\end{english}
\end{pycell}
\\ \hline
$2$ &
\begin{pycell}
\begin{english}
\begin{minted}{python}
for i in range(5, 10):
    print(i)
\end{minted}
\end{english}
\end{pycell}
& $B$ &
\begin{pycell}
\begin{english}
\begin{minted}{python}
x = 2
while x < 13:
    print(x)
    x = x + 2
\end{minted}
\end{english}
\end{pycell}
\\ \hline
$3$ &
\begin{pycell}
\begin{english}
\begin{minted}{python}
i = 1
while i <= 6:
    print(i * 2)
    i = i + 1
\end{minted}
\end{english}
\end{pycell}
& $C$ &
\begin{pycell}
\begin{english}
\begin{minted}{python}
for i in range(5, 31, 5):
    print(i)
\end{minted}
\end{english}
\end{pycell}
\\ \hline
$4$ &
\begin{pycell}
\begin{english}
\begin{minted}{python}
i = 16
while i > 2:
    print(i)
    i = i - 3
\end{minted}
\end{english}
\end{pycell}
& $D$ &
\begin{pycell}
\begin{english}
\begin{minted}{python}
item = 4
while item < 9:
    print(item + 1)
    item = item + 1
\end{minted}
\end{english}
\end{pycell}
\\ \hline
\end{longtable}}

{\renewcommand{\arraystretch}{1.3}
\begin{longtable}{|C{4.5cm}|C{6cm}|}
\hline
\textbf{رقم الحلقة في العمود أ} & \textbf{الحلقة المكافئة من العمود ب ($A$ / $B$ / $C$ / $D$)} \\
\hline
$1$ & \ifwithsols $\mathbf{C}$ (مضاعفات الـ $5$ من $5$ إلى $30$) \fi \\ \hline
$2$ & \ifwithsols $\mathbf{D}$ (الأعداد من $5$ إلى $9$) \fi \\ \hline
$3$ & \ifwithsols $\mathbf{B}$ (الأعداد الزوجية من $2$ إلى $12$) \fi \\ \hline
$4$ & \ifwithsols $\mathbf{A}$ (الأعداد $16, 13, 10, 7, 4$) \fi \\ \hline
\end{longtable}}
\filbreak

% ---------------------------------------------------------------
\item ($12$ علامة) في لعبة \textbf{"التنين الذهبي"}، يحصل اللاعب على رتبة قوّة وفقًا للنقاط التي جمعها في اللعبة. رتب القوّة الممكنة هي:
محارب مبتدئ (\textenglish{Beginner Warrior})، فارس فضّي (\textenglish{Silver Knight})، فارس ذهبي (\textenglish{Gold Knight})، وساحر خارق (\textenglish{Super Wizard}).
اللاعب الذي لديه أقل من $50$ نقطة يُعرَّف برتبة "مبتدئ" (\textenglish{Beginner}).
تُحدَّد رتبة القوّة وفقًا للسيناريو الذي أمامكم:

\begin{boxCode}
\begin{english}
\begin{minted}{python}
points = int(input("Enter Point: "))
if points < 50:
    print("Beginner")
elif points < 200:
    print("Beginner Warrior")
elif points < 500:
    print("Silver Knight")
elif points < 1000:
    print("Gold Knight")
else:
    print("Super Wizard")
\end{minted}
\end{english}
\end{boxCode}

\begin{enumerate}[label=\alph*.]
\item ($4$ علامات) معطى الجدول التالي. اكتبوا ما هي رتبة القوّة للاعب لكل واحدة من قيم النقاط في الجدول:

{\renewcommand{\arraystretch}{1.4}
\begin{longtable}{|C{3cm}|C{6cm}|}
\hline
\textbf{النقاط} & \textbf{رتبة القوّة} \\
\hline
$25$ & \ifwithsols مبتدئ (\textenglish{Beginner}) \fi \\ \hline
$200$ & \ifwithsols فارس فضّي (\textenglish{Silver Knight}) \fi \\ \hline
$555$ & \ifwithsols فارس ذهبي (\textenglish{Gold Knight}) \fi \\ \hline
$1050$ & \ifwithsols ساحر خارق (\textenglish{Super Wizard}) \fi \\ \hline
\end{longtable}}
\filbreak

\item ($8$ علامات) معطى الجدول التالي. أكملوا "رتبة القوّة" وفقًا للبرنامج (رتب القوّة الممكنة: مبتدئ، محارب مبتدئ، فارس فضّي، فارس ذهبي، ساحر خارق):

{\renewcommand{\arraystretch}{1.4}
\begin{longtable}{|C{1.5cm}|>{\raggedright\arraybackslash}p{7cm}|C{5cm}|}
\hline
\textbf{البند} & \centering\textbf{مجال النقاط} & \textbf{رتبة القوّة} \\
\hline
مثال & \textenglish{\texttt{points < 50}} & مبتدئ (\textenglish{Beginner}) \\ \hline
$1$ & \textenglish{\texttt{points < 1000 and points >= 500}} & \ifwithsols فارس ذهبي (\textenglish{Gold Knight}) \fi \\ \hline
$2$ & \textenglish{\texttt{points >= 200 and points < 500}} & \ifwithsols فارس فضّي (\textenglish{Silver Knight}) \fi \\ \hline
$3$ & \textenglish{\texttt{points < 200 and points >= 50}} & \ifwithsols محارب مبتدئ (\textenglish{Beginner Warrior}) \fi \\ \hline
$4$ & \textenglish{\texttt{points >= 1000}} & \ifwithsols ساحر خارق (\textenglish{Super Wizard}) \fi \\ \hline
\end{longtable}}
\filbreak
\end{enumerate}
\filbreak

% ---------------------------------------------------------------
\item ($12$ علامة) في لعبة \textbf{"صائد الكنوز"}، تُعتبر النقاط "جيدة" إذا كانت \textbf{فوق $100$ وحتى $300$ (يشمل $300$)}.
ينفّذ اللاعب $10$ مهام مختلفة، وفي كل مهمة يحصل على نقاط وفقًا لنجاحه.
يجب على اللعبة أن تعدّ وتطبع كم مرة حقّق اللاعب نقاطًا "جيدة".
أمامكم ثلاثة مقاطع برنامج. حدّدوا لكل مقطع برنامج هل ينفّذ المطلوب في السؤال. إذا لا، اشرحوا كيف يمكن تصحيح الكود لكي ينفّذ المطلوب.

\begin{enumerate}[label=\textbf{مقطع \alph*:}, wide=0pt, itemsep=1.5em]

\item .
\begin{boxCode}
\begin{english}
\begin{minted}{python}
count = 0
for i in range(10):
    score = int(input("Enter score: "))
    if score > 100 and score <= 300:
        count = count + 1
print(count)
\end{minted}
\end{english}
\end{boxCode}
\ifwithsols
\begin{boxSolution}
\textbf{نعم، ينفّذ المطلوب.} الحلقة تتكرّر $10$ مرات، تستقبل النقاط في كل مهمة، وتفحص بشكل صحيح إذا كانت أكبر من $100$ وأصغر من أو تساوي $300$.
\end{boxSolution}
\else
\vspace{0.5cm}
\textbf{هل ينفّذ الكود المطلوب؟} \hrulefill \\
 \\
\textbf{طريقة التصحيح (إن لزم):} \hrulefill
\fi
\filbreak

\item .
\begin{boxCode}
\begin{english}
\begin{minted}{python}
count = 0
for i in range(10):
    score = int(input("Enter score: "))
    if score > 100 or score <= 300:
        count = count + 1
print(count)
\end{minted}
\end{english}
\end{boxCode}
\ifwithsols
\begin{boxSolution}
\textbf{لا ينفّذ المطلوب.} لأنّه استعمل العامل \textenglish{or} بدل \textenglish{and}، فيتحقّق الشرط لكل عدد مهما كانت قيمته. \\
\textbf{التصحيح:} تغيير \textenglish{or} إلى \textenglish{and} في سطر الشرط: \textenglish{\texttt{if score > 100 and score <= 300:}}.
\end{boxSolution}
\else
\vspace{0.5cm}
\textbf{هل ينفّذ الكود المطلوب؟} \hrulefill \\
 \\
\textbf{طريقة التصحيح (إن لزم):} \hrulefill
\fi
\filbreak

\item .
\begin{boxCode}
\begin{english}
\begin{minted}{python}
count = 0
for i in range(1, 10, 1):
    score = int(input("Enter score: "))
    if score > 100:
        if score <= 300:
            count = count + 1
print(count)
\end{minted}
\end{english}
\end{boxCode}
\ifwithsols
\begin{boxSolution}
\textbf{لا ينفّذ المطلوب.} لأنّ الحلقة \textenglish{\texttt{range(1, 10, 1)}} تتكرّر $9$ مرات فقط (من $1$ إلى $9$) بدل $10$ مرات. \\
\textbf{التصحيح:} تغيير مدى الحلقة إلى \textenglish{\texttt{range(10)}} أو \textenglish{\texttt{range(1, 11, 1)}}.
\end{boxSolution}
\else
\vspace{0.5cm}
\textbf{هل ينفّذ الكود المطلوب؟} \hrulefill \\
 \\
\textbf{طريقة التصحيح (إن لزم):} \hrulefill
\fi
\filbreak

\end{enumerate}
\filbreak

\end{enumerate}

\clearpage
\section*{الفصل الثاني — اختياري ($40$ علامة)}
في هذا الفصل $4$ أسئلة. يجب الإجابة عن \textbf{سؤالين} من بين الأسئلة $6$–$9$ (لكل سؤال — $20$ علامة).

\begin{enumerate}[resume, label=\textbf{السؤال \arabic*:}, wide=0pt, itemsep=2.2em]

% ---------------------------------------------------------------
\item ($20$ علامة) أمامكم مقطع برنامج يستقبل عددين صحيحين موجبين وينفّذ حسابًا معيّنًا:

\begin{boxCode}
\begin{english}
\begin{minted}{python}
num1 = int(input('Please enter number '))
num2 = int(input('Please enter number '))
while num1 >= num2:
    num1 = num1 - num2
print(num1)
\end{minted}
\end{english}
\end{boxCode}

\begin{enumerate}[label=\alph*.]
\item أكملوا الجدول التالي لقيم \textenglish{num1} و \textenglish{num2} المعطاة:

{\begin{english}
\renewcommand{\arraystretch}{1.4}
\begin{longtable}{|C{3cm}|C{3cm}|C{3.5cm}|}
\hline
\textbf{num1} & \textbf{num2} & \textarabic{\textbf{المُخرج}} \\
\hline
$11$ & $4$ & \ifwithsols $\mathbf{3}$ \fi \\ \hline
$9$ & $3$ & \ifwithsols $\mathbf{0}$ \fi \\ \hline
$5$ & $2$ & \ifwithsols $\mathbf{1}$ \fi \\ \hline
$2$ & $5$ & \ifwithsols $\mathbf{2}$ \fi \\ \hline
\end{longtable}
\end{english}}
\filbreak

\item ماذا ينفّذ مقطع البرنامج؟ أحيطوا بدائرة الإجابة الصحيحة:
\begin{enumerate}[label=\arabic*.]
\item يحسب ويطبع حاصل ضرب \textenglish{num1} بـ \textenglish{num2}.
\item يحسب ويطبع الفرق بين \textenglish{num1} و \textenglish{num2}.
\item يحسب ويطبع باقي القسمة (المودولو) من قسمة \textenglish{num1} على \textenglish{num2}.
\item يحسب ويطبع القوّة لـ \textenglish{num1} بالأسّ \textenglish{num2}.
\end{enumerate}

\ifwithsols
\begin{boxSolution}
الإجابة الصحيحة هي \textbf{$3$}: يحسب ويطبع باقي القسمة (المودولو) من قسمة \textenglish{num1} على \textenglish{num2} عن طريق الطرح المتكرّر لـ \textenglish{num2} من \textenglish{num1} طالما أنّ $\text{num1} \geq \text{num2}$.
\end{boxSolution}
\fi
\filbreak
\end{enumerate}
\filbreak

% ---------------------------------------------------------------
\item ($20$ علامة)
\begin{enumerate}[label=\alph*.]
\item ($8$ علامات) أمامكم رقعة رسم. طول وعرض كل مربّع (خلية) في الرقعة هو $20$ خطوة. وفقًا للكود التالي، ارسموا على الرقعة ما سيرسمه السلحفاة (علماً أنّه يبدأ من النقطة $(0,0)$ متّجهًا لليمين):

\begin{boxCode}
\begin{english}
\begin{minted}{python}
import turtle
wn = turtle.Screen()
player = turtle.Turtle()
player.shape("turtle")
player.goto(0, 0)
for i in range(4):
    player.forward(60)
    player.left(90)
player.forward(20)
for i in range(4):
    player.forward(40)
    player.left(90)
turtle.mainloop()
\end{minted}
\end{english}
\end{boxCode}

\begin{center}
\begin{english}
\begin{tikzpicture}[scale=0.65]
\draw[step=1cm, gray, dashed] (-4,-4) grid (4,4);
\filldraw (0,0) circle (2pt) node[below left] {\small\textbf{(0,0)}};
\ifwithsols
\draw[ultra thick, blue] (0,0) rectangle (3,3);
\draw[ultra thick, red] (1,0) rectangle (3,2);
\fi
\end{tikzpicture}
\end{english}
\end{center}

\ifwithsols
\begin{boxSolution}
الحلقة الأولى ترسم مربّعًا كبيرًا طول ضلعه $60$ خطوة ($3$ خلايا) من $(0,0)$ إلى $(3,3)$. بعد ذلك يتقدّم السلحفاة $20$ خطوة (خلية واحدة لليمين) إلى $(1,0)$، وترسم الحلقة الثانية مربّعًا داخليًا طول ضلعه $40$ خطوة (خليتان) يشترك مع المربّع الكبير في الضلع السفلي والضلع الأيمن.
\end{boxSolution}
\fi
\filbreak

\item ($12$ علامة — $3 \times 4$ علامات) أمامكم ثلاثة مقاطع برنامج. كل مقطع برنامج يرسم شكلًا مختلفًا. لائموا بين كل كود والشكل الذي يرسمه في الجدول أدناه (افترضوا وجود أسطر التهيئة المشتركة في بداية كل كود):

{\renewcommand{\arraystretch}{1.3}
\begin{longtable}{|C{0.9cm}|p{7.2cm}|C{0.9cm}|C{5.2cm}|}
\hline
\multicolumn{2}{|c|}{\textbf{العمود أ (الكود)}} & \multicolumn{2}{c|}{\textbf{العمود ب (الشكل)}} \\
\hline
$1$ &
\begin{pycell}
\begin{english}
\begin{minted}{python}
for index in range(3):
    player.forward(60)
    player.left(90)
turtle.mainloop()
\end{minted}
\end{english}
\end{pycell}
& $A$ &
\begin{english}
\vspace{0.1cm}
\begin{tikzpicture}[scale=0.55]
\draw[thick] (0,0.7) circle (0.7cm);
\draw[thick] (0,1.4) circle (1.4cm);
\filldraw (0,0) circle (3pt);
\end{tikzpicture}
\end{english}
\\ \hline
$2$ &
\begin{pycell}
\begin{english}
\begin{minted}{python}
for index in range(3):
    player.left(90)
    player.forward(30)
    player.right(90)
    player.forward(30)
turtle.mainloop()
\end{minted}
\end{english}
\end{pycell}
& $B$ &
\begin{english}
\vspace{0.1cm}
\begin{tikzpicture}[scale=0.55]
\draw[thick] (0,0) -- (2,0) -- (2,2) -- (0,2);
\filldraw (0,2) circle (3pt);
\end{tikzpicture}
\end{english}
\\ \hline
$3$ &
\begin{pycell}
\begin{english}
\begin{minted}{python}
for i in range(1, 3):
    player.circle(10 * i)
turtle.mainloop()
\end{minted}
\end{english}
\end{pycell}
& $C$ &
\begin{english}
\vspace{0.1cm}
\begin{tikzpicture}[scale=0.5]
\draw[thick] (0,0) -- (0,1) -- (1,1) -- (1,2) -- (2,2) -- (2,3) -- (3,3);
\filldraw (3,3) circle (3pt);
\end{tikzpicture}
\end{english}
\\ \hline
\end{longtable}}

{\renewcommand{\arraystretch}{1.3}
\begin{longtable}{|C{4.5cm}|C{4.5cm}|}
\hline
\textbf{العمود أ ($1$ / $2$ / $3$)} & \textbf{العمود ب ($A$ / $B$ / $C$)} \\
\hline
$1$ & \ifwithsols $\mathbf{B}$ ($3$ أضلاع من مربّع) \fi \\ \hline
$2$ & \ifwithsols $\mathbf{C}$ (درج مكوّن من $3$ درجات) \fi \\ \hline
$3$ & \ifwithsols $\mathbf{A}$ (دائرتان بنصفي قطرين $10$ و $20$) \fi \\ \hline
\end{longtable}}
\filbreak
\end{enumerate}
\filbreak

% ---------------------------------------------------------------
\item ($20$ علامة) \textbf{النص "المضاعَف زوجيًا"} هو نص جديد يُبنى من نص موجود، بحيث نضيف الرمز \textenglish{\texttt{'\^{}'}} بعد كل حرف يقع في مكان (مؤشر) \textbf{زوجي} في النص الأصلي.
\begin{itemize}
\item \textbf{مثال $1$:} للنص \textenglish{\texttt{"Monkey"}}، ينتج النص المضاعَف زوجيًا التالي: \textenglish{\texttt{"M\^{}on\^{}ke\^{}y"}}.
\item \textbf{مثال $2$:} للنص \textenglish{\texttt{"Tiger"}}، ينتج النص المضاعَف زوجيًا التالي: \textenglish{\texttt{"T\^{}ig\^{}er\^{}"}}.
\end{itemize}
العملية التي أمامكم تستقبل نصًا \textenglish{msg} من المستخدم، وتعرض نصًا "مضاعفًا زوجيًا". أكملوا الفراغات $(1)$–$(5)$ في العملية:

\begin{boxCode}
\begin{english}
\begin{minted}{python}
def duplicate_even_chars():
    msg = (1)_________("Please enter word: ")
    result = ""
    for i in range((2)_________):
        result = result + (3)_________
        if i % 2 == 0:
            result = result + (4)_________
    print((5)_________)
\end{minted}
\end{english}
\end{boxCode}

\ifwithsols
\begin{boxSolution}
\begin{enumerate}[label=(\arabic*)]
\item \textenglish{\texttt{input}}
\item \textenglish{\texttt{len(msg)}}
\item \textenglish{\texttt{msg[i]}}
\item \textenglish{\texttt{"\^{}"}}
\item \textenglish{\texttt{result}}
\end{enumerate}
\end{boxSolution}
\fi
\filbreak

% ---------------------------------------------------------------
\item ($20$ علامة) \textbf{العدد "الموقِظ"} هو عدد صحيح موجب يكون فيه مجموع جميع منازله \textbf{أكبر بـ $3$ أضعاف أو أكثر} من عدد المنازل في العدد.
\begin{itemize}
\item \textbf{مثال $1$:} العدد $129$ هو عدد "موقِظ" لأنّ مجموع منازله ($1 + 2 + 9 = 12$) أكبر بـ $3$ أضعاف أو أكثر من عدد منازله ($3$). الحساب: $12 / 3 \geq 3$.
\item \textbf{مثال $2$:} العدد $5$ هو عدد "موقِظ" لأنّ مجموع منازله ($5$) أكبر بـ $3$ أضعاف أو أكثر من عدد منازله ($1$). الحساب: $5 / 1 \geq 3$.
\item \textbf{مثال $3$:} العدد $124$ ليس عددًا "موقِظًا" لأنّ مجموع منازله ($1 + 2 + 4 = 7$) ليس أكبر بـ $3$ أضعاف أو أكثر من عدد منازله ($3$). الحساب: $7 / 3 < 3$.
\end{itemize}

\begin{enumerate}[label=\alph*.]
\item ($4$ علامات) حسب التعريف أعلاه، هل $4158$ هو عدد "موقِظ"؟ إذا نعم، اعرضوا حسابًا رياضيًا (كما في المثال) يُظهر ذلك؛ وإلّا اشرحوا.
\ifwithsols
\begin{boxSolution}
\textbf{نعم}، العدد $4158$ هو عدد "موقِظ" لأنّ مجموع منازله هو $4 + 1 + 5 + 8 = 18$، وعدد منازله هو $4$، ويتحقّق:
$$18 / 4 = 4.5 \geq 3 \quad (\text{أو } 18 \geq 3 \cdot 4 = 12)$$
\end{boxSolution}
\else
\\ \underline{\hspace{14cm}}
\fi
\filbreak

\item ($16$ علامة) أمامكم عملية تستقبل عددًا من المستخدم وتفحص هل العدد "موقِظ". أكملوا الفراغات $(1)$–$(5)$ في العملية:

\begin{boxCode}
\begin{english}
\begin{minted}{python}
def awake_number():
    sum = (1)________
    count = 0
    num = int(input("Please enter a number:"))
    while num > 0:
        sum = (2)_________ + num % 10
        count = count + (3)___________
        num = num // 10
    if (4)________ >= 3 * (5)__________:
        print("The number is awakening")
    else:
        print("The number is not awakening")
\end{minted}
\end{english}
\end{boxCode}

\ifwithsols
\begin{boxSolution}
\begin{enumerate}[label=(\arabic*)]
\item $0$
\item \textenglish{\texttt{sum}}
\item $1$
\item \textenglish{\texttt{sum}}
\item \textenglish{\texttt{count}}
\end{enumerate}
\end{boxSolution}
\fi
\filbreak
\end{enumerate}
\filbreak

\end{enumerate}

\vspace{1cm}
\begin{flushleft}
أرجو لكم وقتًا ممتعًا.

الأستاذ محمود اغبارية.
\end{flushleft}

\end{document}
